\section{Call-by-push-value}

\subsection{Morphisms of CBPV Models}

There are two relevant notions of \emph{morphism} of CBPV models:
\emph{strict} and \emph{lax}.
% morphisms of CBPV models
Given call-by-push-value models
$\mathcal M_1 = (\mathcal V_1, \mathcal E_1, \arr_1, U_1, F_1)$ and
$\mathcal M_2 = (\mathcal V_2, \mathcal E_2, \arr_2, U_2, F_2)$,
A \emph{strict} morphism $G$ from $M_1$ to $M_2$ is given by a pair of functors
$G_{\mathcal{V}}: V_1 \to V_2$ and $G_{\mathcal{E}} : E_1 \to E_2$
that strictly presere the constructors:
\begin{enumerate}
  \item $G_{\mathcal{E}} \circ F_1 = F_2 \circ G_{\mathcal{V}}$
  \item $G_{\mathcal{V}} \circ U_1 = U_2 \circ G_{\mathcal{E}}$
  \item $G_{\mathcal{E}}(A \arr_1 B) = G_{\mathcal{V}}(A) \arr_2 G_{\mathcal{E}}(B)$
  \item $G_{\mathcal{V}}(A_1 \times_1 A_2) = G_{\mathcal{V}}(A_1) \times_2 G_{\mathcal{V}}(A_2)$
  \item $G_{\mathcal{V}}(1_1) = 1_2$
\end{enumerate}
As well as strictly preserving the corresponding universal morphisms
and coherence isomorphisms.

A lax morphism instead preserves these only up to transformation
\begin{enumerate}
  \item $G_{\mathcal{E}} \circ F_1 \Rightarrow F_2 \circ G_{\mathcal{V}}$
  \item $G_{\mathcal{V}} \circ U_1 \Rightarrow U_2 \circ G_{\mathcal{E}}$
  \item $G_{\mathcal{E}}(A \arr_1 B) \Rightarrow G_{\mathcal{V}}(A) \arr_2 G_{\mathcal{E}}(B)$
  \item $G_{\mathcal{V}}(A_1 \times_1 A_2) \Rightarrow G_{\mathcal{V}}(A_1) \times_2 G_{\mathcal{V}}(A_2)$
  \item $G_{\mathcal{V}}(1_1) \Rightarrow 1_2$
\end{enumerate}
Additionally a lax morphism should have a relationship between these
transformations and the universal morphisms, but we will only consider
lax morphisms of thin categories, where such conditions hold
trivially.

\subsection{Kleisli Actions of CBPV Type Constructors}

In CBPV models, all the type constructors are interpreted as functors:
\begin{enumerate}
\item $\to : \op\calV \times \calE \to \calE$
\item $\times : \calV \times \calV \to \calV$
\item $F : \calV \to \calE$
\item $U : \calE \to \calV$
\end{enumerate}
That is, they all have functorial actions on \emph{pure} morphisms of
value types and \emph{linear} morphisms of computation types.
%
We use these functorial actions extensively in the construction of
casts and their corresponding perturbations. But when defining
downcasts of value types and upcasts of computation types, we
additionally need a second functorial action of these categories:
functoriality in \emph{impure} morphisms of value types and
\emph{non-linear} morphisms of computation types. These notions of
morphism are given by the \emph{Kleisli} categories $\calVk$ and
$\calEk$ which have value types and computation types as objects but
morphisms are defined as
\[ \calVk(A,A') = \calE(F A, FA')\]
\[ \calEk(B,B') = \calV(U B, U B')\]
with composition given by composition in $\calE/\calV$.  That is we
need to define a second functorial action, that agrees with the above
on objects for these Kleisli categories:
\begin{enumerate}
\item $\tok : \op\calVk \square \calEk \to \calEk$
\item $\timesk : \calVk \square \calVk \to \calVk$
\item $\Fk : \calVk \to \calEk$
\item $\Uk : \calEk \to \calVk$
\end{enumerate}
Note that rather than the product of categories we use the ``funny
tensor product'' $\square$. This is because the action on
impure/non-linear morphisms for $\tok/\timesk$ do not satisfy ``joint
functoriality'' but instead only ``separate functoriality'', meaning
we give rather than an action on morphisms in both categories
simultaneously instead an action on each argument categories morphisms
with the object in the other category fixed. The existence of these
functorial actions for $\tok$ and $\timesk$ is reliant on the
\emph{strength} of the adjunction. We describe them using the internal
language of CBPV in order to more easily verify their
existence/functoriality:
\begin{enumerate}
\item For $\tok$ we define for $\phi : \calE(F A,F A')$ and $B \in \calE$ the morphism $\phi \tok B : \calV(U(A' \to B),U(A\to B))$ as
  \[ t:U(A'\to B) \vdash \phi \tok B = \{ \lambda x. x' \leftarrow \phi\,[\ret x]; ! t x'\} : U(A \to B) \]
  and for $A \in \calV$ and $f : \calV(UB,UB')$ we define $A \tok f : \calV(U(A \to B),U(A\to B'))$ as
  \[ t : U(A \to B) \vdash A \tok f = \{ \lambda x. !f[\{ ! t x \}]\} \]
\item For $\timesk$ we define for $\phi : \calE(F A_1,FA_2)$ and $A' \in \calV$ the morphism $\phi \timesk A_2$ as
  \[ \bullet : F(A_1\times A_2) \vdash \phi \timesk A_2 = (x_1,x_2) \leftarrow \bullet; x_1' \leftarrow \phi[\ret x_1]; \ret (x_1',x_2) : F(A_1'\times A_2)\]
  and $A_1 \timesk \phi$ is defined symmetrically.
\item For $\Uk$ we need to define for $f : \calV(UB,UB')$ a morphism $\Uk f : \calE(FUB,FUB')$. This is simply given by the functorial action of $F$: $\Uk f = F(f)$
\item Similarly $\Fk \phi = U\phi$
\end{enumerate}

Functoriality in each argument is easily established, meaning for
example for the function type is functorial in each argument:
\begin{enumerate}
\item $(\phi \circ \phi') \tok B = (\phi' \tok B) \circ (\phi \tok B)$
\item $\id \tok B = \id$
\item $A \tok (f \circ f') = (A \tok f) \circ (A \tok f)$
\item $A \tok \id = \id$
\end{enumerate}

Finally, note that all of these constructions lift to squares in a
double CBPV model since the squares themselves form a CBPV model and
the projection functions preserve CBPV structure. For instance, given a square
$\alpha : \phi \ltdyn_{F c_o}^{F c_i} \phi'$ and a horizontal morphism $d : B \rel B'$ of appropriate type, we get a square
\[ \alpha \tok d : \phi \tok B \ltdyn_{U(c_o \to d)}^{U(c_i \to d)} \phi' \tok B' \]